\documentclass[12pt]{article}
\usepackage{amsmath, amssymb, geometry}
\geometry{margin=1in}
\setlength{\parindent}{0pt}
\setlength{\parskip}{1em}

\title{Solutions -- Practice Problem Set 6}
\author{ECON 101 -- Macroeconomics}
\date{Fall 2025}

\begin{document}
\maketitle

\section*{Problem 1}

The Bank of Canada observes that inflation is 3.8\%, above the 2\% target midpoint. It chooses \textbf{contractionary monetary policy}.

\subsection*{(a) Implementing contractionary policy with the key policy rate}

\textbf{Step 1: The policy instrument.}
\begin{itemize}
    \item The Bank of Canada's main policy instrument is the \textbf{target overnight interest rate} (the key policy rate).
    \item Contractionary (tight) monetary policy means the Bank \emph{raises} this target rate.
\end{itemize}

\textbf{Step 2: Operating band and settlement balances.}
\begin{itemize}
    \item The Bank sets an operating band around the target overnight rate (for example, $\pm 0.25$ percentage points).
    \item The \textbf{Bank Rate} is at the top of the band and the \textbf{deposit rate} is at the bottom.
    \item To push the market overnight rate \emph{up} toward its higher target, the Bank of Canada typically:
    \begin{itemize}
        \item reduces the supply of settlement balances in the LVTS / payments system,
        \item or drains liquidity so that banks are more likely to borrow at higher rates.
    \end{itemize}
\end{itemize}

\newpage
\textbf{Step 3: Open-market operations.}
\begin{itemize}
    \item The Bank carries out \textbf{open-market sales} of government securities.
    \item Commercial banks and dealers pay for these securities, which reduces their cash/reserves.
    \item With fewer reserves, the supply of overnight funds falls, putting \emph{upward} pressure on the overnight rate.
    \item The market overnight rate rises and aligns with the higher target rate.
\end{itemize}

\textbf{Conclusion:}
\begin{itemize}
    \item Higher overnight interest rate $\Rightarrow$ tighter monetary conditions.
    \item This is the first step in contractionary monetary policy aimed at reducing inflation from 3.8\% toward 2\%.
\end{itemize}

\subsection*{(b) Monetary transmission mechanism (step-by-step)}

Starting from a higher overnight interest rate:

\textbf{Step 1: From overnight rate to other interest rates.}
\begin{itemize}
    \item Higher overnight rate increases short-term money market rates.
    \item Banks' funding costs rise, so they raise prime lending rates, mortgage rates, and other loan rates.
    \item Long-term bond yields may also rise as expectations adjust.
\end{itemize}

\textbf{Step 2: Effect on consumption (\(C\)).}
\begin{itemize}
    \item Higher interest rates make borrowing for consumer durables (cars, appliances) and housing more expensive.
    \item Households reduce spending on credit-sensitive items.
    \item Wealth effects: higher rates may reduce asset prices (e.g., bonds, equities, housing), making households feel poorer and spend less.
    \item Overall, planned consumption spending falls.
\end{itemize}


\newpage
\textbf{Step 3: Effect on investment (\(I\)).}
\begin{itemize}
    \item Higher interest rates increase the cost of borrowing for firms.
    \item Some investment projects that were profitable at lower interest rates become unprofitable.
    \item Firms cut back on new capital spending (machinery, buildings, technology).
    \item Planned investment spending falls.
\end{itemize}

\textbf{Step 4: Effect on exchange rate and net exports (\(NX\)).}
\begin{itemize}
    \item Higher Canadian interest rates attract foreign capital (higher return on CAD assets).
    \item Demand for Canadian dollars rises, causing the \textbf{Canadian dollar to appreciate}.
    \item An appreciation makes Canadian exports more expensive and imports cheaper.
    \item Net exports (NX) decline.
\end{itemize}

\textbf{Step 5: Effect on aggregate demand (AD).}
\begin{itemize}
    \item Lower $C$, lower $I$, and lower $NX$ all reduce total planned expenditure.
    \item In the AD--AS framework, this reduction appears as a \textbf{leftward shift of the AD curve}.
\end{itemize}

\textbf{Step 6: Effect on real GDP and inflation.}
\begin{itemize}
    \item In the short run, with an upward-sloping SRAS, a leftward shift of AD reduces real GDP and the price level.
    \item Lower real GDP (below potential) increases unemployment.
    \item Excess capacity and weaker demand put downward pressure on inflation.
    \item Over time, inflation falls back toward the 2\% target.
\end{itemize}

\subsection*{(c) AD--AS diagram and returning inflation to target}

\textbf{Step 1: Initial equilibrium.}
\begin{itemize}
    \item Initially, the economy is at a point where AD intersects SRAS and LRAS at an inflation rate of 3.8\%.
    \item LRAS is vertical at potential GDP $Y^*$.
\end{itemize}

\textbf{Step 2: Contractionary monetary policy.}
\begin{itemize}
    \item Higher policy interest rate shifts AD leftward (from AD\(_0\) to AD\(_1\)).
    \item In the short run, the new equilibrium is at the intersection of AD\(_1\) and SRAS.
    \item At this point:
    \begin{itemize}
        \item Real GDP \(Y_1 < Y^*\),
        \item Price level and inflation are lower than before.
    \end{itemize}
\end{itemize}

\textbf{Step 3: Medium-run adjustment.}
\begin{itemize}
    \item The output gap (Y below potential) eventually puts downward pressure on wages and other input costs.
    \item SRAS shifts to the right over time, moving the economy back toward potential GDP $Y^*$ at a lower inflation rate.
    \item The final outcome is output near $Y^*$ and inflation close to the 2\% target midpoint.
\end{itemize}

\textbf{Summary:}
\begin{itemize}
    \item Contractionary monetary policy: AD shifts left.
    \item Short run: lower output, higher unemployment, lower inflation.
    \item Medium run: SRAS shifts right, output returns to potential, inflation stabilizes near target.
\end{itemize}



\section*{Problem 2}

We are given the reaction function:
\[
i = 1 + 1.5(\pi - 2) + 0.5(Y - Y^*),
\]
where $i$ is the target overnight interest rate in percent, $\pi$ is inflation, and $(Y - Y^*)$ is the output gap in percent.

Data:
\[
\pi = 3.4\%, \qquad Y - Y^* = +1.5\%.
\]

\subsection*{(a) Compute the recommended overnight interest rate}

\textbf{Step 1: Compute the inflation deviation from target.}
\[
\pi - 2 = 3.4 - 2.0 = 1.4 \quad (\text{percentage points above target}).
\]

\textbf{Step 2: Compute the inflation term $1.5(\pi - 2)$.}
\[
1.5 \times 1.4 = 2.1.
\]

\textbf{Step 3: Compute the output gap term $0.5(Y - Y^*)$.}
\[
Y - Y^* = +1.5 \Rightarrow 0.5 \times 1.5 = 0.75.
\]

\textbf{Step 4: Plug values into the reaction function.}
\[
i = 1 + 1.5(\pi - 2) + 0.5(Y - Y^*) 
= 1 + 2.1 + 0.75 = 3.85.
\]

\textbf{Answer:}
\[
\boxed{i = 3.85\% \text{ (target overnight interest rate)}}
\]

\subsection*{(b) Is this expansionary or contractionary? Explain.}

\textbf{Step 1: Compare to a neutral or lower previous rate.}
\begin{itemize}
    \item The rule indicates a relatively \emph{high} policy rate (3.85\%) because:
    \begin{itemize}
        \item inflation is above target (3.4\% vs 2\%),
        \item and the output gap is positive (economy above potential).
    \end{itemize}
\end{itemize}

\textbf{Step 2: Interpretation.}
\begin{itemize}
    \item Raising the policy rate to 3.85\% will \emph{lean against} inflationary pressures and an overheating economy.
    \item This is a \textbf{contractionary} (tightening) monetary policy stance:
    \begin{itemize}
        \item it is intended to slow down aggregate demand,
        \item reduce the positive output gap,
        \item and bring inflation back toward 2\%.
    \end{itemize}
\end{itemize}

\textbf{Conclusion:}
\[
\boxed{\text{The policy stance implied by } i = 3.85\% \text{ is contractionary.}}
\]

\subsection*{(c) Effects on lending, long-term rates, AD, and inflation}

\textbf{Step 1: Bank lending and short-term rates.}
\begin{itemize}
    \item A higher overnight rate increases banks' cost of short-term funds.
    \item Banks typically respond by raising their prime lending rates and other short-term loan rates.
    \item As borrowing becomes more expensive, households and firms reduce their demand for loans.
    \item Bank lending slows.
\end{itemize}

\textbf{Step 2: Long-term interest rates.}
\begin{itemize}
    \item Expectations of tighter monetary policy raise yields on longer-term bonds.
    \item Mortgage rates and corporate bond yields may rise.
    \item This further dampens spending on housing and business investment.
\end{itemize}

\newpage

\textbf{Step 3: Aggregate demand.}
\begin{itemize}
    \item Higher interest rates reduce:
    \begin{itemize}
        \item \textbf{Consumption} of durable goods (cars, appliances, housing-related spending).
        \item \textbf{Investment} in new capital (machines, buildings, technology).
    \end{itemize}
    \item The Canadian dollar may appreciate, reducing net exports.
    \item Combined, $C$, $I$, and $NX$ fall, causing a leftward shift in \textbf{aggregate demand (AD)}.
\end{itemize}

\textbf{Step 4: Inflation in the medium run.}
\begin{itemize}
    \item With AD lower, real GDP growth slows and the positive output gap shrinks or disappears.
    \item Lower demand pressure reduces upward pressure on wages and prices.
    \item Over time, this moderates inflation, moving it from 3.4\% back toward the 2\% target.
\end{itemize}

\textbf{Summary:}
\begin{itemize}
    \item Higher policy rate $\Rightarrow$ higher borrowing costs and reduced credit growth.
    \item Lower $C$, $I$, and $NX$ $\Rightarrow$ lower AD.
    \item Lower AD $\Rightarrow$ lower inflation in the medium run.
\end{itemize}

\hrule
\vspace{1em}

\begin{center}
\textit{End of Solutions}
\end{center}

\end{document}
